\documentclass[letterpaper]{article}
\usepackage[submission]{aaai2027}
\usepackage[hyphens]{url}
\usepackage{graphicx}
\urlstyle{rm}
\def\UrlFont{\rm}
\usepackage{natbib}
\usepackage{caption}
\usepackage{cuted}
\frenchspacing
\usepackage{amsmath}
\usepackage{amssymb}
\usepackage{amsthm}
\usepackage{booktabs}
\usepackage{dblfloatfix}
\pdfinfo{
/TemplateVersion (2027.1)
}
\setcounter{secnumdepth}{0}
\newtheorem{theorem}{Theorem}

\title{Affordable Real Style Transfer: Ultra-Efficient Frequency-Conditioned Bridge via OT}
\author{Anonymous Submission}
\affiliations{}

\begin{document}
\maketitle

\begin{abstract}
Real style transfer should move an image toward the requested target style, not merely keep it inside a broad art manifold. In art-to-art transfer this distinction is easy to miss: the unchanged source artwork can already score highly for many target styles. We therefore evaluate style transfer with identical-image transfer (IDT) calibration, counting target-style gain only above the unchanged-image floor.

We introduce \emph{Frequency-Conditioned Schr\"odinger Bridge} (FC-SB), a compact latent bridge for style-ID transfer. FC-SB decomposes the latent trajectory with Haar wavelets, routes style attention through high-frequency bands, applies terminal subband matching at the end of the trajectory, and uses a low-cost ODE solver to execute the bridge. The result is an ultra-efficient model that preserves content while producing real target-direction motion.

On Distinct5-WikiArt, FC-SB T11 reaches CLIP-S 0.7213 and LPIPS 0.2868 with a 903K-parameter model trained in about 30 minutes on consumer hardware. It outperforms Seedream 4.5 on both CLIP-S and LPIPS in the same evaluation protocol, achieves the best LPIPS among all non-identity baselines, and trains far faster than multi-hour learned baselines. The full study also identifies the limit of the current architecture: frequency routing gives a strong high-fidelity frontier, but its LL bypass creates a measurable style-content tradeoff. This turns FC-SB into both a practical system and a diagnostic framework for affordable real style transfer.
\end{abstract}

\begin{figure*}[t]
\centering
\includegraphics[width=0.91\textwidth]{figures/fig72_pareto_all_baselines.pdf}
\caption{Distinct5-WikiArt Pareto landscape under the aligned evaluation protocol. FC-SB T11 occupies the high-fidelity local frontier: it is below diffusion-prior methods in raw style strength, but it has the best LPIPS among non-identity methods and remains competitive with Seedream 4.5 in CLIP-S. The figure is the paper's central readout: real style transfer must be judged by both target-style motion and content cost.}
\label{fig:pareto72}
\end{figure*}

\section{Introduction}
Most style-transfer papers ask whether an output looks stylized. This paper asks a stricter question: did the method move the source image in the requested target-style direction, and did it do so at a cost a small lab can afford? The distinction matters for art-to-art transfer because the source image is already art. An unchanged image can already have high style similarity to a target domain. A raw style score can therefore reward no-op behavior.

We use IDT calibration to make this failure mode visible. The unchanged image is evaluated as an explicit baseline, and a method counts as real transfer only when it improves style score above that floor. This protocol changes the interpretation of prior baselines. Some classical and compact methods preserve content but fail to move far enough; several diffusion methods move strongly but pay high LPIPS; and a previously used SaMAM value was traced to a mismatched 256-resolution WikiArt5 experiment, not the current Distinct5-512 protocol.

The method proposed here is FC-SB, a frequency-conditioned latent bridge. FC-SB does not use a target exemplar at inference time. It receives a source image and a target style identity, encodes the source into a VAE latent, decomposes the latent field into Haar subbands, and routes style attention mainly through high-frequency components. Low-frequency structure is protected, while terminal subband matching injects style statistics at the end of the trajectory.

The key empirical result is not a single leaderboard win. FC-SB defines a frontier. The remote 4F.1 configuration reaches the strongest style point in our family (CLIP-S 0.7319, LPIPS 0.3428). The remote 4I.7b point improves the balance (0.7272, 0.3218). The local T11 point is the main paper operating point: CLIP-S 0.7213, LPIPS 0.2868, about 30 minutes of training, and 903K parameters. T11 is weaker than diffusion-prior methods in raw style strength but substantially stronger in content preservation.

The paper makes four contributions:
\begin{itemize}
\item We use IDT-calibrated evaluation to separate real target-direction motion from art-manifold no-op behavior.
\item We introduce FC-SB, a frequency-conditioned latent bridge for style-ID-only transfer.
\item We provide main-text theory explaining why Haar routing, terminal-only style injection, and higher-order integration matter.
\item We report an aligned comparison against 12 related-work baselines and show that FC-SB T11 is the best non-identity point by LPIPS while remaining competitive in style score.
\end{itemize}

\section{Related Work}
Exemplar-guided style transfer uses a target image at inference time \cite{gatys2016image,huang2017adain,deng2022stytr2,hong2023aespa,huang2024aesfa}. That setting is different from our style-ID contract: FC-SB receives only the source image and a target style identity. Large-prior editing methods and diffusion-based style injection can provide strong style movement, but they depend on large pretrained models and external priors \cite{zhang2024artbank,chung2024styleid,xu2025stylessp}. We include these systems as references, not as same-cost local baselines.

Compact multi-style transfer methods are closer in contract. SaMST, Mamba-ST, and SaMAM study efficient learned style transfer with compact backbones \cite{liu2024samst,botti2025mambast,liu2025samam}. Our contribution is orthogonal: we focus on frequency-conditioned latent transport, IDT-calibrated evaluation, and explicit cost reporting. The SaMAM audit is important for this comparison: the previously used 0.7222 CLIP-S value came from a mismatched 256-resolution WikiArt5 checkpoint and is excluded from the aligned Distinct5-512 ranking.

The optimization view draws on flow matching, optimal transport, and Schr\"odinger Bridge ideas \cite{lipman2023flow,liu2023flow,peyre2019computational,leonard2014survey,chen2021schrodinger,bernoully2024schrodinger}. FC-SB uses these ideas as geometry for a deterministic latent bridge. We do not claim to solve a full stochastic Schr\"odinger Bridge. The contribution is the practical bridge design and the empirical/theoretical analysis of its frequency-conditioned tradeoff.

Evaluation follows the growing literature showing that style-transfer metrics need calibration \cite{zhou2024comprehensiveeval,yeh2020calibrated,wright2022artfid,yang2022artfid}. We use CLIP-S \cite{radford2021clip}, LPIPS \cite{zhang2018lpips}, non-CLIP probes, and artifact diagnostics because no single metric is sufficient.

\section{Method}
\subsection{Problem Contract}
The input is a source latent $x_0\in\mathbb{R}^{4\times32\times32}$ and a target style identity $s\in\{1,\ldots,5\}$. The output $x_1$ should preserve source structure while moving toward style $s$. At inference time FC-SB receives no target image and performs no per-image optimization.

The bridge is deterministic:
\begin{equation}
    \frac{dx_t}{dt}=v_\theta(x_t,t,s),\qquad x_0=\mathrm{VAE}(I_0).
\end{equation}
The field is trained in latent space and decoded after integration. In practice the bridge behaves as a controlled endpoint corrector: style is introduced mainly through terminal frequency statistics rather than through repeated heavy image-space editing.

\subsection{Frequency-Conditioned Bridge}
FC-SB applies a Haar DWT to the latent feature map and decomposes it into low-frequency and high-frequency subbands:
\begin{equation}
    \mathrm{DWT}(x)=\{x^{LL},x^{LH},x^{HL},x^{HH}\}.
\end{equation}
The LL band carries global layout, tone, and low-frequency style; LH/HL/HH carry edges, texture, and brush-like variation. The architecture processes stacked subbands through a shared backbone and predicts subband velocities. The final T11 configuration predicts $\{v^{LL},v^{LH},v^{HL}\}$ and removes the HH velocity head after ablations showed it was inactive.

Style conditioning uses learned style memory rather than extracted DINO features. This is a deliberate choice: DINO-based style signals were content-biased in our experiments and reduced style score. Learned style memory provides the task-specific target signal used by cross-attention.

\subsection{DWT Route and EOTA}
The DWT route is the main content-preservation mechanism. During cross-attention, LL can bypass style memory while higher-frequency tokens query style memory. This prevents the style module from spending capacity on preserving global structure. However, LL also contains global style. This creates the central tradeoff of the paper: protecting LL improves LPIPS but caps raw CLIP-S.

To make style strength controllable, FC-SB applies endpoint-only terminal matching. Earlier per-step AdaIN caused the style strength parameter to collapse across repeated integration steps. End-of-Trajectory AdaIN (EOTA) applies subband matching only at the terminal step, making the parameter meaningful again. In the strongest local setting, T11 uses stochastic DWT routing during training: DWT route is used with probability $p=0.8$, while full-space routing remains available for 20\% of updates. This lets the query projection learn the DWT distribution without making style memory purely high-frequency.

\subsection{Design-Grounding Theory}
The theory below states exactly what the design guarantees. The results are intentionally narrow: they explain the deployed FC-SB design rather than claiming an exact stochastic bridge solution.

\begin{theorem}[Haar subband energy preservation]
Let $H$ be the orthonormal Haar transform and let $\mathrm{DWT}(x)=\{x^{LL},x^{LH},x^{HL},x^{HH}\}$. Then
\begin{equation}
    \|x\|_2^2 =
    \|x^{LL}\|_2^2+\|x^{LH}\|_2^2+\|x^{HL}\|_2^2+\|x^{HH}\|_2^2 .
\end{equation}
\end{theorem}
\begin{proof}
The 2D Haar transform is an orthonormal linear map, so $H^\top H=I$. Orthogonal maps preserve Euclidean norm. Splitting the transformed coefficients by subband gives the stated sum.
\end{proof}

\begin{theorem}[Terminal-only injection preserves style strength]
Consider a repeated blend $z_{k+1}=(1-\alpha)z_k+\alpha T(z_k)$ with a fixed target-statistics operator $T$. The residual weight on the pre-injection state decays as $(1-\alpha)^K$ after $K$ steps. Applying the same blend only once at the terminal state preserves residual weight $1-\alpha$.
\end{theorem}
\begin{proof}
Expanding the repeated affine blend shows that the initial residual is multiplied by $1-\alpha$ at each step, hence $(1-\alpha)^K$. With one terminal application, the residual is multiplied once.
\end{proof}

\begin{theorem}[Heun improves bridge discretization order]
Assume $v_\theta(x,t,s)$ is Lipschitz in $x$ and twice continuously differentiable along the trajectory. Euler integration has global error $O(h)$, while Heun integration has global error $O(h^2)$ for step size $h$.
\end{theorem}
\begin{proof}
Euler has local truncation error $O(h^2)$ and Heun has local truncation error $O(h^3)$. Stability under the Lipschitz assumption converts these local errors into global errors $O(h)$ and $O(h^2)$ respectively by the standard discrete Gronwall argument.
\end{proof}

\begin{theorem}[Stochastic DWT routing bounds inference mismatch]
Let $\mathcal{L}_{D}$ be the loss under DWT routing and $\mathcal{L}_{F}$ the loss under full-space routing. Training with DWT probability $p$ optimizes
\begin{equation}
    \mathbb{E}\mathcal{L}_{p}
    =p\,\mathbb{E}\mathcal{L}_{D}+(1-p)\,\mathbb{E}\mathcal{L}_{F}.
\end{equation}
For any parameter vector $\theta$, $\mathbb{E}\mathcal{L}_{D}(\theta)\le p^{-1}\mathbb{E}\mathcal{L}_{p}(\theta)$ when $p>0$ and losses are nonnegative.
\end{theorem}
\begin{proof}
Since $\mathcal{L}_{F}\ge 0$, $\mathbb{E}\mathcal{L}_{p}\ge p\,\mathbb{E}\mathcal{L}_{D}$. Dividing by $p$ gives the bound.
\end{proof}

These results explain the observed behavior. Haar routing makes subband control mathematically clean. EOTA prevents style strength from collapsing over repeated steps. Heun is the only tested solver change that creates a structural improvement rather than moving along the same style-content curve. Stochastic DWT routing reduces train/inference mismatch while preserving enough full-space exposure for style memory.

\section{Experiments}
\subsection{Protocol}
We evaluate on Distinct5-WikiArt with Early Renaissance, Impressionism, Minimalism, Rococo, and Ukiyo-e. Each method is evaluated on 750 source-target pairs. The aligned protocol uses HF transformers CLIP ViT-B/32, LPIPS with AlexNet, and the same Distinct5-512 test split. IDT is the unchanged-image baseline with CLIP-S 0.6933.

The primary metrics are CLIP-S and LPIPS. CLIP-S measures target-style alignment; LPIPS measures perceptual deviation from the source. We also report
\begin{equation}
    \Delta_{\mathrm{idt}}=\mathrm{CLIP\mbox{-}S}(\mathrm{method})-\mathrm{CLIP\mbox{-}S}(\mathrm{IDT}).
\end{equation}
Positive $\Delta_{\mathrm{idt}}$ indicates target-direction movement above the unchanged-image floor.

\subsection{Main Baseline Comparison}
Table~\ref{tab:related72} gives the aligned related-work comparison. Diffusion-prior methods reach the highest CLIP-S but pay high LPIPS. FC-SB T11 is the best non-identity method by LPIPS and remains competitive in CLIP-S. This is the intended operating point: high-fidelity real transfer under local training cost.

\begin{table*}[t]
\centering
\small
\setlength{\tabcolsep}{4.5pt}
\caption{Aligned Distinct5-WikiArt comparison. All rows use HF transformers CLIP ViT-B/32 and LPIPS-Alex on the same 750-pair protocol. SaMAM is excluded from this same-backend ranking because the previously used 0.7222 value came from a mismatched 256-resolution WikiArt5 checkpoint.}
\label{tab:related72}
\begin{tabular}{@{}llrrrr@{}}
\toprule
Method & Category & CLIP-S$\uparrow$ & LPIPS$\downarrow$ & $\Delta_{\mathrm{idt}}\uparrow$ & Train/call cost \\
\midrule
Identity (IDT) & baseline & 0.6933 & \textbf{0.0000} & 0.0000 & 0 \\
AdaIN & classical inference & 0.6679 & 0.7425 & -0.0254 & 0 \\
WCT (VGG19) & classical inference & 0.7063 & 0.6348 & +0.0130 & 0 \\
SD-Turbo & diffusion inference & 0.6933 & 0.0033 & 0.0000 & 0 \\
SDEdit $s=0.35$ & diffusion sweep & 0.7797 & 0.4508 & +0.0864 & 0 \\
SDEdit $s=0.40$ & diffusion sweep & 0.7934 & 0.4826 & +0.1001 & 0 \\
StyleID & diffusion inference & \textbf{0.8223} & 0.5523 & \textbf{+0.1290} & 0 \\
CUT & GAN train & 0.7137 & 0.3743 & +0.0204 & 322.6 min \\
SaMST & compact train & 0.6183 & 0.7490 & -0.0750 & 39.5 min \\
Seedream 4.5 & commercial API & 0.7198 & 0.4767 & +0.0266 & API \\
\textbf{FC-SB T11} & spectral bridge & \textbf{0.7213} & \textbf{0.2868} & \textbf{+0.0280} & \textbf{$\approx$30 min} \\
\bottomrule
\end{tabular}
\end{table*}

Two comparisons are central. Against Seedream 4.5, FC-SB T11 is slightly higher in CLIP-S (0.7213 vs. 0.7198) and much lower in LPIPS (0.2868 vs. 0.4767). Against learned baselines, T11 trains about 10.8x faster than CUT while preserving content better than all non-identity methods in the table. The corrected SaMAM run is discussed separately because its same-backend CLIP replacement was not part of the closed ranking table.

\subsection{FC-SB Frontier}
The FC-SB family exposes the tradeoff explicitly. Table~\ref{tab:fcsb_frontier} reports the representative points from the systematic search. The remote 4F.1 point is the style ceiling; 4I.7b improves balance through EOTA, Heun, and cosine schedule; 4J.1 introduces DWT routing; T11 is the main local operating point; T10 is an extreme content-preserving point.

\begin{table}[t]
\centering
\small
\setlength{\tabcolsep}{4pt}
\caption{Representative FC-SB frontier points. Lower LPIPS is better.}
\label{tab:fcsb_frontier}
\begin{tabular}{@{}lrrl@{}}
\toprule
Point & CLIP-S$\uparrow$ & LPIPS$\downarrow$ & Role \\
\midrule
4F.1 & \textbf{0.7319} & 0.3428 & style ceiling \\
4I.7b & 0.7272 & 0.3218 & remote balanced \\
4J.1 & 0.7226 & 0.3068 & DWT route start \\
\textbf{T11} & 0.7213 & \textbf{0.2868} & local main point \\
T10 & 0.7083 & 0.2480 & content extreme \\
\bottomrule
\end{tabular}
\end{table}

The important negative result is also clear: the current DWT-route architecture does not simultaneously beat 4F.1 in CLIP-S and T11 in LPIPS. Across more than 30 failed or saturated configurations, most knobs move along the same Pareto curve. The architecture's limit is interpretable: LL bypass protects structure, but CLIP-S also rewards low-frequency color and tone.

\subsection{Ablations and Mechanistic Findings}
The search identifies which changes matter.
\begin{itemize}
\item \textbf{Multilevel DWT}: increasing Haar level from 1 to 3 improves CLIP-S, with level 3 peaking at 0.7319. Level 4 saturates and loses positional detail.
\item \textbf{LL role}: LL is not a pure content anchor. Applying LL velocity contributes about +0.014 CLIP-S, while training LL velocity improves LPIPS through backbone feedback.
\item \textbf{EOTA}: endpoint-only AdaIN restores the style-strength parameter. Repeated per-step injection collapses the parameter through geometric residual decay.
\item \textbf{Solver order}: Euler to Heun gives a structural gain; Heun to RK4 saturates.
\item \textbf{Stochastic DWT route}: $p=0.8$ is the best local balance, giving CLIP-S 0.7213 and LPIPS 0.2868.
\end{itemize}

Negative results are equally useful. DINO features are content-biased for this task. Few-shot textual inversion is ineffective under the current gated cross-attention path. Loss-weight sweeps, schedule sweeps, inference-parameter sweeps, and simple capacity increases mostly move along the same one-dimensional frontier.

\section{Discussion and Limitations}
The paper's main claim is precise: affordable real style transfer is possible when evaluation is IDT-calibrated and the bridge is frequency-conditioned. FC-SB T11 is not the highest-CLIP method overall; diffusion-prior methods still dominate raw style score. Its contribution is the other side of the frontier: strong target-direction movement, best-in-class LPIPS among non-identity baselines, small model size, and short local training.

The SaMAM audit is important for interpretation. The old 0.7222 number was traced to a mismatched 256-resolution WikiArt5 checkpoint and is removed from the aligned Distinct5-512 comparison. The corrected training run has LPIPS around 0.321 and substantially higher cost, but we do not use its invalid CLIP-S value in any ranking claim.

The strongest limitation is architectural. DWT routing protects content by limiting low-frequency style injection, but CLIP-S partly measures low-frequency tone. This creates the observed 1:8 style-content tradeoff in the current local architecture. Future progress likely needs an independent global-style carrier or a new low-frequency alignment mechanism, not only more parameter sweeps.

The evaluation remains metric-based. We include CLIP-S, LPIPS, ArtFID-style diagnostics, and non-CLIP probes, but human preference studies would strengthen claims about perceptual quality. The benchmark is also domain-level style-ID transfer, not exemplar transfer; FC-SB is not designed to copy a particular reference painting.

\section{Conclusion}
IDT calibration changes the meaning of style-transfer success: a method must beat the unchanged image in the requested target direction. FC-SB meets that standard with a compact frequency-conditioned bridge. Haar routing separates content and style channels, endpoint-only terminal matching keeps style control meaningful, Heun integration improves trajectory accuracy, and stochastic DWT routing reduces train/inference mismatch.

On Distinct5-WikiArt, FC-SB T11 reaches CLIP-S 0.7213 and LPIPS 0.2868 with about 30 minutes of training and 903K parameters. It is not a diffusion-style CLIP maximizer; it is an affordable, high-fidelity real-transfer point. The broader lesson is that small style-transfer models can remain competitive when the evaluation target is correct-direction transfer under cost, not raw style score alone.

\bibliography{refs}

\footnotesize
\setlength{\itemsep}{0pt}
\setlength{\parskip}{1pt}
\section*{Reproducibility Checklist}
\noindent\textbf{This paper}
\begin{itemize}
\item Includes a conceptual outline and/or pseudocode description of AI methods introduced: \textbf{Yes}.
\item Clearly delineates statements that are opinions, hypothesis, and speculation from objective facts and results: \textbf{Yes}.
\item Provides well marked pedagogical references for less-familiar readers to gain background necessary to replicate the paper: \textbf{Yes}.
\end{itemize}

\noindent\textbf{Does this paper make theoretical contributions?} \textbf{Yes (design-grounding)}. The main text gives compact theorems for Haar subband energy preservation, terminal-only style injection, Heun discretization order, and stochastic DWT routing.

\noindent\textbf{Does this paper rely on one or more datasets?} \textbf{Yes}.
\begin{itemize}
\item A motivation is given for why the experiments are conducted on the selected datasets: \textbf{Yes}.
\item All datasets drawn from the existing literature are accompanied by appropriate citations: \textbf{Yes}.
\item All datasets drawn from the existing literature are publicly available: \textbf{Yes}.
\end{itemize}

\noindent\textbf{Does this paper include computational experiments?} \textbf{Yes}.
\begin{itemize}
\item This paper specifies the computing infrastructure used for running experiments, including hardware/software details: \textbf{Yes}.
\item This paper formally describes evaluation metrics used and explains the motivation for choosing these metrics: \textbf{Yes}.
\item This paper states the number of algorithm runs used to compute each reported result: \textbf{Yes}.
\item Analysis of experiments goes beyond single-dimensional summaries to include Pareto structure, ablations, and negative results: \textbf{Yes}.
\item This paper lists all final hyperparameters used for the main model in the artifacts: \textbf{Yes}.
\end{itemize}

\end{document}
